\subNsection{1}{KI-Reflexion}{Kristian Popov}

Mein Schwerpunkt lag während des Projekts auf den Themen \emph{Endliche Automaten}
und \emph{Reguläre Sprachen}.
Dabei habe ich Gemini 3.1 Pro, die kostenlose Version von ChatGPT sowie Claude Opus
4.8 und später Opus 5 ausprobiert.

Die Modelle waren dabei unterschiedlich gut.
Manche erklärten theoretische Zusammenhänge besser, während andere bei der
Implementierung oder Fehlersuche hilfreicher waren.
Überwiegend arbeitete ich mit Claude, da die Vorschläge bei Programmieraufgaben
häufiger direkt brauchbar waren und ich dadurch weniger nacharbeiten musste.

Im Laufe des Projekts ist mir außerdem aufgefallen, dass die Modelle unterschiedlich
viel Aufwand in die Planung einer Antwort stecken.
Eine höhere Planungsstufe bedeutet mehr verbrauchte Tokens und längere Wartezeiten,
führt aber nicht automatisch zu besseren Ergebnissen.
Für kleine Änderungen reichte meist eine niedrige Stufe aus, während komplexere
Aufgaben mit mehreren Dateien von einer höheren Stufe profitierten.
Deshalb habe ich je nach Aufgabe entschieden, welche Einstellung sinnvoll war.

\subNsection{2}{KI als Co-Developer}{Kristian Popov}

Die KI habe ich direkt in die Entwicklungsumgebung eingebunden.
Dadurch konnte sie auf Projektdateien zugreifen, Änderungen vorschlagen und diese
direkt im Projekt überprüfen.

Besonders hilfreich war sie bei wiederkehrenden Aufgaben wie Refactorings,
Boilerplate-Code, Unit-Tests, Internationalisierung und dem Beheben von
TypeScript-Warnungen.
Bei der Automatenlogik und beim Parsen von Grammatiken konnte sie außerdem schnell
passende Testfälle ableiten.

Bei Architekturentscheidungen ließ ich mir zunächst mehrere Lösungsansätze mit ihren
Vor- und Nachteilen erklären und traf die Entscheidung anschließend selbst.
Jede vorgeschlagene Änderung habe ich überprüft und getestet.

Unterstützt wurde diese Arbeitsweise durch Erweiterungen aus dem offiziellen
Plugin-Verzeichnis \texttt{claude-\allowbreak plugins-\allowbreak official}.
Aktiv waren mehrere, darunter \texttt{typescript-lsp} für Typinformationen über das
\gls{lsp}.
Für meine Themen am wichtigsten waren zwei:

\begin{description}[leftmargin=*,style=nextline,itemsep=0.75ex]
\item[\texttt{playwright}]
ermöglicht der KI, die Anwendung im Browser auszuführen und aus Sicht des Benutzers
zu testen.
Dadurch konnten beispielsweise Oberflächen, Navigation und Benutzerinteraktionen
überprüft werden.

\item[\texttt{chrome-devtools-mcp}]
ermöglicht über das \gls{mcp} den Zugriff auf Konsolenausgaben und das Verhalten der
laufenden Anwendung.
Dadurch konnte ich bei Fehlern genauer untersuchen, was zur Laufzeit tatsächlich
passiert.
\end{description}

Zusätzlich habe ich das Projekt einmalig mit Graphify~\cite{shamsi2026graphify}
analysiert und einen Wissensgraphen erzeugt.
Dadurch konnte die KI bei späteren Aufgaben gezielt auf relevante Dateien, Klassen
und Funktionen zugreifen, ohne das gesamte Projekt erneut einlesen zu müssen.

\subNsection{2}{Automatisierte Aufgabengenerierung}{Kristian Popov}

Für die Übungen wurden neben der Aufgabenstellung auch eine Referenzlösung, ein
Startprogramm und Beispielwörter benötigt.
Die KI nutzte ich vor allem, um Varianten bestehender Aufgaben zu erzeugen,
beispielsweise mit anderem Alphabet, anderen Wortlängen oder unterschiedlichen
Schwierigkeitsgraden.

Ob eine Aufgabe tatsächlich korrekt war, habe ich anschließend mit der automatischen
Übungskontrolle des Simulators gegen die Referenzlösung geprüft
(siehe Abschnitt~\ref{4_2_EA}).
Nur geprüfte Aufgaben wurden in den Aufgabenkatalog aufgenommen
(siehe Abbildung~\ref{fig:ki_aufgaben}).

\begin{figure}[htbp]
	\centering
	\begin{tikzpicture}[
		font=\footnotesize,
		box/.style={draw, rounded corners=2pt, align=center, inner sep=5pt, minimum height=1.4cm},
		flow/.style={-{Stealth[length=2.2mm,width=1.8mm]}, semithick},
		lbl/.style={font=\scriptsize\itshape, inner sep=2pt}
	]
		\node[box, text width=2.6cm, fill=black!5]  (a) at (0,0)    {Bestehende Aufgabe\\als Vorlage};
		\node[box, text width=3.4cm, fill=black!5]  (b) at (4.1,0)  {KI erzeugt Variante\\(Alphabet, Wortlänge, Schwierigkeit)};
		\node[box, text width=3.2cm, fill=black!12] (c) at (8.6,0)  {Übungskontrolle prüft gegen die Referenzlösung};
		\node[box, text width=2.6cm, fill=black!5]  (d) at (13.2,0) {Aufnahme in den Aufgabenkatalog};

		\draw[flow] (a) -- (b);
		\draw[flow] (b) -- (c);
		\draw[flow] (c) -- (d) node[lbl, midway, above=2pt, fill=white] {bestanden};
		\draw[flow] (c.south) -- ++(0,-1.0) coordinate (fb) -| (b.south);
		\node[lbl, below=1mm] at ($(fb)!0.5!(fb -| b.south)$)
			{Abweichung: Variante wird verworfen};
	\end{tikzpicture}
	\caption[Erzeugung und Prüfung von Aufgabenvarianten]{%
		Erzeugung und Prüfung von Aufgabenvarianten. Nur Varianten, die sich
		gegenüber der Referenzlösung identisch verhalten, werden in den
		Aufgabenkatalog übernommen.}
	\label{fig:ki_aufgaben}
\end{figure}

\subNsection{2}{Kritische Reflexion}{Kristian Popov}

Die KI lieferte häufig brauchbare, aber nicht immer korrekte Vorschläge.
Teilweise verwendete sie APIs, die in der verwendeten Version nicht mehr existierten,
oder erfand Funktionen und Bibliotheken.
Auch Sonderfälle und Auswirkungen auf andere Teile des Projekts wurden gelegentlich
übersehen.

Ein weiteres Problem war die starke Ähnlichkeit zwischen den Themen
\emph{Endliche Automaten} und \emph{Reguläre Sprachen}.
Statt bestehende Funktionen zu abstrahieren, erzeugte die KI teilweise doppelte
Implementierungen.
Änderungen, die mehrere Themen gleichzeitig betrafen, musste ich deshalb häufig in
kleinere Schritte aufteilen.

Insgesamt sparte mir die KI besonders bei wiederkehrenden Aufgaben und beim Ausbau
bestehender Funktionen viel Zeit.
Bei theoretischen Konzepten und größeren Architekturänderungen war jedoch weiterhin
meine eigene Prüfung entscheidend.
Ein Orchestrator wie Sakana Fugu könnte künftig die Auswahl des passenden Modells
übernehmen~\cite{sakana2026fugu}; da das System erst kurz vor Abschluss meiner Arbeit
erschien, konnte ich es nicht mehr einsetzen.

Für mich war die KI damit vor allem ein Werkzeug zur Unterstützung.
Die fachlichen Entscheidungen und die Verantwortung für die Korrektheit der
Implementierung lagen jederzeit bei mir.
